\scanpage{4}
$E$ faisceau loc. libre de rang $r+1$ sur $S$

$P=\mathrm{P}(E)$ fibré projectif associé, on a une suite exacte canonique
\[
0\longrightarrow\Omega^1_{P/S}\longrightarrow g^*(\underline E)(-1)\longrightarrow\underline O_P\longrightarrow0
\]
[qui donne un isom can
\[
\Omega^r_{P/S}\xrightarrow{\sim}g^*\bigl(\bigwedge_{\underline O_S}^{r+1}E\bigr)(-r-1)\tag*{(1)}
\]
\hfill$]$

Si \unsure{donc} $X$ est un sous-préschéma de $P$ \unsure{simple|lisse} sur $S$, de dim relative \inserted{$n=$}$r-p$
associé au faisceau \deleted{loc. libre}\inserted{d'idéaux} $J$, on a la suite exacte
\[
0\longrightarrow J/J^2\longrightarrow\Omega^1_{P/S}\otimes\underline O_X\longrightarrow\Omega^1_{X/S}\longrightarrow0
\]
d'où l'isomorphisme
\[
\Omega^r_{P/S}\otimes\underline O_X\xleftarrow{\sim}\Omega^{\deleted{\unsure{r-p}}\unsure{n}}_{X/S}\otimes\bigwedge^p_{\underline O_X}(J/J^2)\tag*{(2)}
\]
d'où en composant (1) et (2), désignant par $f:X\to S$
le morphisme structural
\[
\deleted{\Omega^r}f^*\bigl(\bigwedge_{\underline O_S}^{r+1}\underline E\bigr)(-r-1)\xleftarrow{\sim}\Omega^{\deleted{\unsure{r-p}}\unsure{n}}_{X/S}\otimes\bigwedge^p_{\underline O_X}(J/J^2)
\]

Supposons que $X$ soit une intersection complète, définie
par un sous-module \inserted{gradué $M$ \deleted{\unsure{de|\(\subset\)}}} loc. facteur direct de $\mathrm{Symm.}^{*}(\underline E)$,
\deleted{\unsure{qui soit de rang}} loc. libre de rang $p$ [il n'est \deleted{de fait}
déterminé que modulo le carré $\underline I^2$ de l'idéal qu'il
définit ; mais sa structure est connue]. On a
un hom, si $I_a$ est la composante de degré $a$ de $I$
\[
g^*(I_d)(-d)\longrightarrow J
\]
d'où des hom \uline{surjectif}
\[
\coprod g^*(M_d)(-d)\longrightarrow J
\]
induisant un \uline{isom}
\[
\coprod f^*(M_d)(-d)\xrightarrow{\sim}J/J^2\simeq\mathrm{Conormal}_{P/X}
\]
% archive sheet number: 3

